\documentclass[12pt]{article}
\usepackage{amsmath, amssymb}

\title{Lecture Scribe: Inequalities and Tail Bounds}
\author{}
\date{}

\begin{document}

\maketitle

\section{Tail of a Random Variable and Tail Probability}

\subsection{Definition}

The tail of a random variable $X$ is defined as:
\[
P(X > x)
\]

This represents the probability that the random variable takes a value greater than a threshold $x$.

\section{Motivation for Studying Tail Events}

\subsection{Why do we care about tails?}

\begin{itemize}
    \item Jobs delayed more than 24 hours
    \item Packets dropped due to full buffer
\end{itemize}

These represent rare but critical events.

\subsection{Key Idea}

\begin{itemize}
    \item Mean $\rightarrow$ average behavior (easy to analyze)
    \item Tail $\rightarrow$ rare extreme events (hard to analyze)
\end{itemize}

\subsection{Why are tails hard to compute?}

\begin{itemize}
    \item Requires summing many small probabilities
\end{itemize}

\section{Exact Computation of Tail Probabilities}

\subsection{Example 1: Binomial Case}

Let:
\[
X \sim \text{Binomial}(n, p)
\]

We want:
\[
P(X \ge k)
\]

Exact computation:
\[
P(X \ge k) = \sum_{i=k}^{n} \binom{n}{i} p^i (1-p)^{n-i}
\]

\subsection{Example 2: Poisson Case}

Let:
\[
X \sim \text{Poisson}(\lambda)
\]

Then:
\[
P(X \ge k) = \sum_{i=k}^{\infty} e^{-\lambda} \frac{\lambda^i}{i!}
\]

\subsection{Observation}

\begin{itemize}
    \item These are long, complicated summations
    \item Hence exact computation is difficult
\end{itemize}

\newpage
\section{Idea of Tail Bounds}

\subsection{Motivation}

Instead of computing exact values:

\begin{itemize}
    \item Find an upper bound
    \item Easier to compute
    \item Still provides a guarantee
\end{itemize}

\[
P(X \ge k) \le \text{(simple expression)}
\]

\subsection{Key Idea}

Approximate safely instead of computing exactly.

\section{Running Example}

\begin{itemize}
    \item Flip a fair coin $n$ times
    \item Let $X$ = number of heads
\end{itemize}

Distribution:
\[
X \sim \text{Binomial}(n, 1/2)
\]

Mean:
\[
E[X] = \frac{n}{2}
\]

Tail event:
\[
P\left(X \ge \frac{3n}{4}\right)
\]

Interpretation:

This corresponds to getting unusually many heads.
\newpage
\section{Markov's Inequality}

\subsection{Assumptions}

\[
X \ge 0
\]

\subsection{Statement}

\[
P(X \ge a) \le \frac{E[X]}{a}
\]

\subsection{Key Idea}

\begin{itemize}
    \item Uses only the mean
    \item Large values are rare
\end{itemize}

\subsection{Proof}

\begin{enumerate}
    \item Write expectation:
    \[
    E[X] = \sum_x x \cdot P(X = x)
    \]
    
    \item Restrict to values where $X \ge a$:
    \[
    E[X] \ge \sum_{x \ge a} x \cdot P(X = x)
    \]
    
    \item Since $x \ge a$:
    \[
    E[X] \ge \sum_{x \ge a} a \cdot P(X = x)
    \]
    
    \item Factor out $a$:
    \[
    E[X] \ge a \cdot P(X \ge a)
    \]
    
    \item Rearranging:
    \[
    P(X \ge a) \le \frac{E[X]}{a}
    \]
\end{enumerate}

\subsection{Example}

If:
\[
E[X] = 10
\]

Then:
\[
P(X \ge 100) \le \frac{10}{100} = 0.1
\]

\section{Chebyshev's Inequality}

\subsection{Assumptions}

Mean $\mu$, variance $\sigma^2$

\subsection{Statement}

\[
P(|X - \mu| \ge k) \le \frac{\sigma^2}{k^2}
\]

\subsection{Derivation}

\begin{enumerate}
    \item Let:
    \[
    Y = (X - \mu)^2 \ge 0
    \]
    
    \item Apply Markov:
    \[
    P(Y \ge k^2) \le \frac{E[Y]}{k^2}
    \]
    
    \item Since:
    \[
    E[Y] = \sigma^2
    \]
    
    \item Hence:
    \[
    P((X - \mu)^2 \ge k^2) \le \frac{\sigma^2}{k^2}
    \]
    
    \item Equivalent to:
    \[
    P(|X - \mu| \ge k) \le \frac{\sigma^2}{k^2}
    \]
\end{enumerate}
\newpage
\section{Chernoff Bound}

\subsection{Goal}

Provide exponentially tighter bounds.

\subsection{Key Idea}

\[
P(X \ge a) = P(e^{tX} \ge e^{ta})
\]

Apply Markov:
\[
P(e^{tX} \ge e^{ta}) \le \frac{E[e^{tX}]}{e^{ta}}
\]

Thus:
\[
P(X \ge a) \le \frac{E[e^{tX}]}{e^{ta}}
\]

\subsection{Optimization}

Choose $t$ to minimize the bound.

\section{Poisson Tail Bounds}

For:
\[
X \sim \text{Poisson}(\lambda)
\]

Exact:
\[
P(X \ge k) = \sum_{i=k}^{\infty} e^{-\lambda} \frac{\lambda^i}{i!}
\]

Instead, use Chernoff bounds.

\section{Jensen's Inequality}

\subsection{Definition}

A function $f$ is convex if:
\[
f(\theta x + (1-\theta)y) \le \theta f(x) + (1-\theta)f(y)
\]

\subsection{Statement}

\[
f(E[X]) \le E[f(X)]
\]

\subsection{Example}

If $f(x) = x^2$:
\[
(E[X])^2 \le E[X^2]
\]

\section{Connections}

\subsection{Core Flow}

\begin{enumerate}
    \item Exact tails are hard
    \item Use inequalities:
    \begin{itemize}
        \item Markov (mean)
        \item Chebyshev (variance)
        \item Chernoff (exponential)
    \end{itemize}
    \item Goal:
    \[
    P(X \ge k) \le \text{simple bound}
    \]
\end{enumerate}

\subsection{Hierarchy}

\begin{itemize}
    \item Markov $\rightarrow$ weakest
    \item Chebyshev $\rightarrow$ stronger
    \item Chernoff $\rightarrow$ strongest
\end{itemize}

\section{Final Takeaways}

\begin{itemize}
    \item Tail probability: $P(X > x)$
    \item Exact computation is difficult
    \item Inequalities provide upper bounds
    \item Markov $\rightarrow$ mean
    \item Chebyshev $\rightarrow$ variance
    \item Chernoff $\rightarrow$ exponential bounds
    \item Jensen connects convexity and expectation
\end{itemize}

\[
\text{Bound rare events instead of computing them exactly}
\]

\end{document}